\chapter{Contrato del modelo de decisión}

El contrato entre los modelos generados y el simulador se mantiene
deliberadamente pequeño. Un modelo válido implementa:

\begin{lstlisting}[language=Python]
def decide(self, perception: dict) -> Action
def update(self, action, reward, new_perception) -> None
def get_state(self) -> dict
\end{lstlisting}

\section{Percepción}

La percepción es un diccionario construido por el entorno. Las claves mínimas
usadas por los modelos actuales son:

\begin{verbatim}
x, y, grid_width, grid_height, step,
resources, last_action_result
\end{verbatim}

El modelo no debe modificar la percepción recibida. Debe leerla, escoger una
acción y esperar a \texttt{update} para cambiar su estado interno.

\section{Acciones}

Cada modelo define una clase \texttt{Action} ligera. Esta decisión evita importar
clases internas del simulador en código generado. El simulador solo necesita que
la acción exponga los campos esperados por el entorno.

\section{Estado}

\texttt{get\_state} devuelve un diccionario serializable con variables internas.
En modelos de aprendizaje por refuerzo debe incluir valores como
\texttt{q\_values}. Esto permite observar el comportamiento antes y después de
cada paso de simulación.
